\documentclass[11pt]{article}

\usepackage[T1]{fontenc}
\usepackage{amsmath,amssymb,amsthm,mathtools}
\usepackage[a4paper,margin=29mm]{geometry}
\usepackage[hidelinks]{hyperref}
\usepackage{microtype}

\newtheorem{theorem}{Theorem}[section]
\newtheorem{lemma}[theorem]{Lemma}
\newtheorem{corollary}[theorem]{Corollary}
\newtheorem{remark}[theorem]{Remark}
\newcommand{\Mort}{\mathsf{M}}
\newcommand{\Z}{\mathbb Z}
\newcommand{\Q}{\mathbb Q}
\newcommand{\eps}{\varepsilon}
\newcommand{\pref}{\preccurlyeq}

\title{Four-Generator Generalized Post Correspondence\
and Mortality of Five $3\times3$ Integer Matrices}
\author{GPT-5.6 Sol\\\normalsize elicited by \texttt{@eternalism\_4eva}}
\date{21 July 2026}

\begin{document}
\maketitle

\begin{abstract}
Neary reduced an undecidable restricted binary-tag halting problem to five-pair
Post correspondence, but the published converse contains a defective terminal-pair
argument.  We prove directly that the four nonterminal pairs satisfy a fixed-right-boundary
word equation exactly when the source tag system halts.  The proof uses an exhaustive
finite-state synchronization argument and a queue-history invariant; it covers arbitrary
prefix-compatible deviations rather than only intended simulations.  This yields
undecidability of binary-target generalized PCP with four source generators.  A fresh-marker
construction also gives a corrected binary five-pair PCP family whose primitive solutions
end in the fifth pair.

Applying the standard ternary matrix representation and one rank-one boundary separator,
we obtain undecidability of mortality for at most five $3\times3$ integer matrices.  Four
generators are nonsingular and upper triangular; the fifth is nonzero of rational rank one.
The resulting minimal known undecidability antichain is
\[
  \Mort_3(5),\quad \Mort_5(4),\quad \Mort_6(3),\quad \Mort_{12}(2).
\]
The Lean~4 development formalizes the four-tile source equivalence, fresh-marker repair, and
exact integer-matrix equivalence without project axioms or admitted statements.
\end{abstract}

\section{Introduction}

For a finite set $\mathcal A$ of square matrices, mortality asks whether the zero matrix
belongs to the multiplicative semigroup $\mathcal A^+$.  Write $\Mort_d(k)$ for the
restriction to sets of at most $k$ rational $d\times d$ matrices.  Restricting inputs to
integer matrices does not alter the decidability status.  Cassaigne, Halava, Harju, and
Nicolas~\cite{CHHN14} established the then-minimal known undecidability antichain
\[
  \Mort_3(6),\qquad \Mort_5(4),\qquad \Mort_9(3),\qquad \Mort_{15}(2),
\]
in their Table~5 and recorded $\Mort_3(5)$ as open.  Later surveys retained this
bound~\cite{BPS21,Dong23}.

The apparent route through Neary's five-pair PCP theorem~\cite{Neary15} has a subtle defect.
Its claimed exclusion of the terminal pair uses a false assertion about adjacent encoded
symbols.  Rote diagnosed this and proposed a longer terminal block~\cite{Rote25}.  That
repair controls intended tag configurations, but its printed argument does not classify all
malformed prefix-comparable PCP residuals.  We avoid both arguments.  Equality of the two
PCP words itself forces the four ordinary tiles into exact deletion-width blocks.  Decoding
those blocks yields a global queue equation, and a generic prefix-cancellation lemma turns
that equation into a lawful tag execution, stopping at the first short queue.  Post-halting
or malformed residuals are therefore irrelevant.

Our source result is the following.  Here $\operatorname{GPCP}(k)$ counts letters in the
source alphabet of the two morphisms, following Nicolas~\cite{Nicolas08}.

\begin{theorem}\label{thm:gpcp}
$\operatorname{GPCP}(4)$ over a binary target alphabet is undecidable.  It remains
undecidable with empty left boundaries, one empty right boundary, and nonempty morphism
images.
\end{theorem}

The matrix consequence uses the established fixed-boundary compiler, whose rank-one form
also appears in earlier forced-endpoint constructions~\cite{HHH07,CHHN14}.

\begin{theorem}\label{thm:main}
$\Mort_3(5)$ is undecidable.  It remains undecidable under the promise that four generators
are nonsingular upper-triangular integer matrices and the fifth is a nonzero integer matrix
of rank one over $\Q$.
\end{theorem}

Neary's one-version 2013 technical report~\cite{Neary13} already claimed the same numerical
statement: its Theorem~3 claimed four-pair PCP undecidable, and Corollary~2 derived
$\Mort_3(5)$.  That premise was not retained in the refereed 2015 successor, which explicitly
left four-pair PCP open and derived only the six-matrix bound.  CHHN, submitted after the
preprint, likewise recorded $\Mort_3(5)$ as unknown, and later literature retained their
bound.  We found no formal withdrawal, published diagnosis of the precise 2013 defect, or
independent proof of its claim.  Thus the present theorem is not the first public statement of
$\Mort_3(5)$; it is a proof of the previously unestablished case and, to our knowledge, the
first valid one.  Any priority claim should remain qualified pending author review.

\section{Restricted tag systems and four word pairs}

Fix an integer $\beta\geq3$ and a word $q\in\{b,c\}^*$.  Consider the deterministic tag
system with deletion number $\beta$, rules
\[
  b\longmapsto b,\qquad c\longmapsto qb,
\]
and initial queue
\[
  I=\operatorname{drop}_{\beta-1}(q)b.
\]
A step deletes the first $\beta$ symbols and appends the rule word selected by the first
deleted symbol.  The system halts when its queue has length less than $\beta$.

Put
\[
  H(b)=10^\beta1,\qquad H(c)=1,\qquad M=10^\beta,
\]
and extend $H$ morphically to words.  Define four pairs, labelled by a rule symbol and an
erase symbol:
\begin{align*}
 R_c&=(1,\;1H(q)10),& R_b&=(10^\beta1,\;110),\\
 D_c&=(1,\;0),&       D_b&=(10^\beta1,\;0).
\end{align*}
For a label word $w$, write $U(w)$ and $V(w)$ for the concatenated upper and lower words.

\begin{theorem}[Four-tile tag compiler]\label{thm:source}
Suppose $\beta-1\mid |q|$ and $|q|\geq\beta-1$.  Then
\[
  \exists w\in\{R_c,R_b,D_c,D_b\}^*:U(w)M=V(w)
  \quad\Longleftrightarrow\quad
  T_{\beta,q}\text{ halts from }I.
\]
\end{theorem}

The proof occupies the remainder of this section.  Its soundness direction is the point
missing from the earlier terminal-tile arguments.

\subsection{Block forcing}

A \emph{stroke} is a word $a_0a_1\cdots a_{\beta-1}\in\{b,c\}^\beta$.  Associate to it the
tile block
\[
  B(a_0\cdots a_{\beta-1})
  =R_{a_0}D_{a_1}\cdots D_{a_{\beta-1}}.
\]
For a stroke history $S=(s_1,\ldots,s_n)$, let $C(S)=s_1\cdots s_n$ be the word consumed
by the strokes and let $P(S)$ be the concatenation of the tag appendants selected by their
first symbols.

\begin{lemma}[Synchronization]\label{lem:synchronization}
If $U(w)M=V(w)$, then $w=B(s_0)B(s_1)\cdots B(s_n)$ for a nonempty stroke history.
\end{lemma}

\begin{proof}
The word $U(w)M$ begins with $1$, every positive run of zeros in it has length exactly
$\beta$, and it ends with a run of $\beta$ zeros.  Scan the equal word $V(w)$ while retaining
only the length of the current zero run.  From the initial state an erase tile is impossible.
Either rule tile moves the run length from $0$ to $1$.  Each erase tile increases it by one;
a further rule tile is compatible only after the count reaches $\beta$; and an erase after
$\beta$ is impossible.  Acceptance at the terminal zero run therefore forces one rule tile
followed by exactly $\beta-1$ erase tiles, repeatedly.  The labels on their upper words give
the corresponding strokes.  The empty history cannot satisfy the original equality.
\end{proof}

This scan is a two-shape automaton with states ``before the first $1$'' and ``$j$ zeros since
the last $1$'', $0\leq j\leq\beta$.  Thus Lemma~\ref{lem:synchronization} is exhaustive over
all label words; it does not presume that a prefix follows the intended tag simulation.

\subsection{Queue histories}

Appending one final $1$ completes the dangling marker: $M1=H(b)$.  Direct expansion of a
block history gives
\begin{align}
 U(B(S))&=H(C(S)),\label{eq:upper-history}\\
 V(B(s_0)B(S'))1
   &=H\bigl(c\,Q(\operatorname{head}s_0)\,b\,P(S')\bigr),
                                                        \label{eq:lower-history}
\end{align}
where $Q(c)=q$ and $Q(b)=\eps$.  Since $H$ is injective, the terminal equation and
Equations~\eqref{eq:upper-history}--\eqref{eq:lower-history} first force
$\operatorname{head}s_0=c$.  If $s_0=c z$ with $|z|=\beta-1$, cancellation then yields
\begin{equation}\label{eq:init-history}
  z\,C(S')b=q\,b\,P(S').
\end{equation}
Both $z$ and $q$ are prefixes of the same word, and $|z|=\beta-1\leq |q|$.  Hence
$z=\operatorname{take}_{\beta-1}(q)$.  Cancelling it from~\eqref{eq:init-history} gives
\begin{equation}\label{eq:queue-certificate}
  C(S')b=I\,P(S').
\end{equation}

\begin{lemma}[History soundness]\label{lem:history}
Let a tag system have deletion width $\beta$.  If a stroke history $S$, initial queue $I$,
and word $F$ with $|F|<\beta$ satisfy
\[
  C(S)F=I\,P(S),
\]
then the tag system halts from $I$.
\end{lemma}

\begin{proof}
Induct on $S$.  For the empty history, $I=F$ is already short.  Otherwise, if $I$ is short,
stop.  If not, the first stroke and $I$ are both prefixes of the common word in the displayed
equality, while the stroke has length $\beta\leq|I|$.  It is therefore the first $\beta$
symbols of $I$.  Delete it, append the rule word selected by its head, cancel the common
prefix in the global equation, and invoke the induction hypothesis.  Notice that the proof
deliberately stops if a short queue occurs before the prescribed end of the history.
\end{proof}

Taking $F=b$ in Lemma~\ref{lem:history} and using~\eqref{eq:queue-certificate} proves the
soundness direction of Theorem~\ref{thm:source}.

For the converse, record the $\beta$ deleted symbols of each lawful tag step as a stroke.
Every resulting queue ends in $b$.  Moreover,
\[
  |I|\equiv1\pmod{\beta-1},\qquad
  |b|\equiv|qb|\equiv1\pmod{\beta-1},
\]
so every reachable queue has length congruent to $1$ modulo $\beta-1$.  A reachable queue
of length less than $\beta$ must consequently be the one-letter queue $b$.  The execution
therefore supplies a history satisfying~\eqref{eq:queue-certificate}.  Prepending the
initial stroke $c\operatorname{take}_{\beta-1}(q)$ and reversing the calculation above
constructs a terminal match.  This completes the proof of Theorem~\ref{thm:source}.

\subsection{The undecidable source family}

The arithmetic hypotheses above do not characterize the outputs of Neary's compiler; they
define a strictly larger envelope on which Theorem~\ref{thm:source} remains true.  We now
verify explicitly that an undecidable compiler-output subfamily lies inside that envelope.

\begin{lemma}[Computable congruent padding]\label{lem:padding}
Let a cyclic-tag instance have program length $p$, input length $n$, and maximum appendant
length $r$.  Put $\beta=10p$ and
\[
 B=\max\{11(n+p)+\beta-2,\ 11p,\ 11r,\ 1\}.
\]
There is a computable choice
\[
 x=\left\lceil\frac{B-1}{\beta-1}\right\rceil,
 \qquad s=x(\beta-1)+1,
\]
for which Neary's Table~2 construction is defined and its simulation proof is unchanged.
In particular, $s\ge B$ and $s\equiv1\pmod{\beta-1}$.
\end{lemma}

\begin{proof}
The ceiling formula gives $s\ge B$ and the congruence immediately.  The variable exponents
in Table~2 are
\[
 s-11(n+p)-\beta+2,\quad s-11p+1,\quad s-11p,\quad
 s-11v+1,\quad s-11v,\quad s-1,
\]
where $v\le r$; all are nonnegative under the displayed bound.  Every specified track
therefore still has length $s$.  Increasing $s$ changes only the displayed padding and full
copies of the program word.  Since that word has length $\beta s$, all object lengths and
the shift residues modulo $\beta$ used in Neary's track-by-track simulation are unchanged.
Thus the verification of Table~2 applies to this computably chosen $s$.
\end{proof}

Neary's Lemma~9~\cite{Neary15}, together with Lemma~\ref{lem:padding}, supplies an
undecidable subfamily of the systems in Theorem~\ref{thm:source}.  His whole $c$-appendant
$qb$ has length $\beta s$.  Consequently
\[
  |q|=\beta[x(\beta-1)+1]-1=(x\beta+1)(\beta-1),
\]
so the hypotheses of Theorem~\ref{thm:source} hold, and his input
$q_{\beta-1}q_\beta\cdots b$ is precisely $I$.  We use Neary's tag-universality lemma here,
but no soundness claim from his subsequent five-pair PCP proof.

\begin{proof}[Proof of Theorem~\ref{thm:gpcp}]
Use the four labels as the GPCP source alphabet; use $U,V$ as the two morphisms; take both
left boundaries empty, and take $M,\eps$ as the right boundaries.  The GPCP equation is
exactly $U(w)M=V(w)$.  The construction is computable and all eight ordinary tile words
are nonempty.  The empty source word is not a solution because $M\ne\eps$.
Thus the same family is an instance of the nonempty-witness convention
$\operatorname{GPCP}^{+}(4)$.  Theorem~\ref{thm:source}, Lemma~\ref{lem:padding}, and
Neary's Lemma~9 now give undecidability.
\end{proof}

\section{A replacement five-pair PCP source}\label{sec:marker}

Although GPCP and mortality need only Theorem~\ref{thm:source}, the same argument gives a
corrected replacement for Neary's defective five-pair instance.  It does not validate either
the original four-one guard or Rote's longer unary guard.  Introduce a fresh symbol $\#$,
lift every ordinary tile word from $\{0,1\}^*$ into $\{0,1,\#\}^*$, and add
\[
  (M\#,\;\#).
\]
Finally return to a binary target alphabet with the fixed-length injective code
\[
  \eta(0)=00,\qquad \eta(1)=01,\qquad \eta(\#)=11.
\]

\begin{theorem}\label{thm:five-pcp}
The resulting binary five-pair PCP instance is solvable exactly when $T_{\beta,q}$ halts.
Every primitive solution contains the fifth pair exactly once, in its final position.
\end{theorem}

\begin{proof}
The fixed-length code preserves equality and prefix comparability at tile boundaries.
Every ordinary upper word ends in $1$, whereas every ordinary lower word ends in $0$;
hence an ordinary-only label word cannot be a PCP solution.  Every solution therefore uses
the fifth pair.  At its first occurrence, prefix comparability requires
\[
  U(w)M\#\quad\text{and}\quad V(w)\#
\]
to be prefix-comparable.  Since $\#$ is fresh, words $x\#$ and $y\#$ are prefix-comparable
only when $x=y$.  Thus $U(w)M=V(w)$, and the prefix through this first terminal tile is
already a PCP solution.  If a nonempty suffix remained, cancellation would make that suffix
another PCP solution, contradicting primitivity.  The converse follows by appending the
terminal pair to any terminal match and applying $\eta$.
\end{proof}

This synchronization replaces the quantitative unary guard proposed by Rote.  In particular,
Theorem~\ref{thm:five-pcp} remains valid even if malformed ordinary prefixes wander through
prefix-comparable residuals after a simulated halt.

\section{The matrix compiler}

Recode binary letters as nonzero ternary digits, $0\mapsto1$ and $1\mapsto2$.  For a digit
word $x$, let $\sigma(x)$ be its base-$3$ value, with $\sigma(\eps)=0$, and define
\[
\Psi(x,y)=
\begin{pmatrix}
1 & \sigma(y) & \sigma(x)-\sigma(y)\\
0 & 3^{|y|} & 3^{|x|}-3^{|y|}\\
0 & 0 & 3^{|x|}
\end{pmatrix}.
\]

\begin{lemma}\label{lem:morphism}
For binary-coded words,
\begin{align*}
 \Psi(xx',yy')&=\Psi(x,y)\Psi(x',y'),\\
 e_1^{\mathsf T}\Psi(x,y)e_3=0&\Longleftrightarrow x=y,\\
 \det\Psi(x,y)&=3^{|x|+|y|}.
\end{align*}
\end{lemma}

\begin{proof}
The first identity follows from
$\sigma(xx')=3^{|x'|}\sigma(x)+\sigma(x')$ by multiplication.  The selected entry is
$\sigma(x)-\sigma(y)$.  Positional notation over the nonzero digits $1,2$ is injective even
between unequal lengths.  The determinant is the product of the diagonal entries.
\end{proof}

For the four ordinary labels put $X_i=\Psi(U_i,V_i)$ after ternary recoding, and put
\[
  X_\star=\Psi(M,\eps),\qquad c=X_\star e_3,
\]
where words in the latter display are likewise recoded.  Define five integer matrices by
\[
  A_i=X_i\quad(1\leq i\leq4),
  \qquad A_\star=X_\star e_3e_1^{\mathsf T}=ce_1^{\mathsf T}.
\]
Each $X_i$ is upper triangular and nonsingular over $\Q$.  The third coordinate of $c$ is
$3^{|M|}$, so $A_\star$ is nonzero and has rank one.

\begin{lemma}[Rank-one boundary compiler]\label{lem:matrix}
The five emitted matrices are mortal if and only if
\[
  \exists w\in\{1,2,3,4\}^*:U(w)M=V(w).
\]
\end{lemma}

\begin{proof}
A terminal match gives
\[
  A_\star X_w A_\star
  =c\bigl(e_1^{\mathsf T}X_wX_\star e_3\bigr)e_1^{\mathsf T}=0
\]
by Lemma~\ref{lem:morphism}.

Conversely, a product without $A_\star$ is nonsingular.  Every product containing $t\geq1$
copies has a unique decomposition
\[
  P_0A_\star P_1A_\star\cdots A_\star P_t,
\]
where each $P_j$ is a possibly empty product of ordinary matrices.  Expanding the outer
products gives
\[
 (P_0c)\left(\prod_{j=1}^{t-1}e_1^{\mathsf T}P_jc\right)
 (e_1^{\mathsf T}P_t).
\]
The exterior column and row are nonzero because $P_0,P_t$ are nonsingular.  For $t=1$ the
outer product is nonzero.  For $t\geq2$ the product vanishes exactly when some internal
scalar vanishes.  Writing the corresponding block as $P_j=X_w$, Lemma~\ref{lem:morphism}
gives
\[
 e_1^{\mathsf T}P_jc=0
 \Longleftrightarrow e_1^{\mathsf T}X_wX_\star e_3=0
 \Longleftrightarrow U(w)M=V(w).
\]
This includes empty exterior or interior blocks and adjacent separators.
\end{proof}

The construction is computable.  Theorem~\ref{thm:source} and Neary's undecidable source
family therefore prove Theorem~\ref{thm:main}.

\section{Undecidability-antichain consequences}

Padding each generator by a zero direct summand preserves mortality.  Cassaigne et al.'s
trade
\[
  \Mort_d(hk+1)\leq_m\Mort_{kd}(h+1)
\]
from their Theorem~4, applied to $\Mort_3(5)$ with $(h,k)=(2,2)$ and $(1,4)$, gives
\[
  \Mort_3(5)\leq_m\Mort_6(3),\qquad
  \Mort_3(5)\leq_m\Mort_{12}(2).
\]
Together with their independent $\Mort_5(4)$ result, the minimal known undecidable antichain
becomes
\[
  \boxed{\Mort_3(5),\quad\Mort_5(4),\quad\Mort_6(3),\quad\Mort_{12}(2)}.
\]

There are parallel consequences in the notation of~\cite{CHHN14}.  Their $Z_d(k)$ is scalar
zero reachability: an instance supplies a row $L$, a column $C$, and at most $k$ matrices,
and asks whether $LYC=0$ for some nonempty product $Y$.  Their $R_d(k)$ asks for a zero in
the upper-right corner.  The proof of their Theorem~1 gives
$\operatorname{GPCP}^{+}(k)\leq_m Z_3(k)$; Theorem~6 preserves the common-first-column
structure, Theorem~7 gives its dimension trade, and Propositions~3--4 convert $Z$ to $R$
and mortality.  Together with their Theorem~3, these reductions give the minimal known
undecidability antichains
\[
  Z_3(4),\ Z_5(3),\ Z_7(2),
  \qquad
  R_3(5),\ R_4(4),\ R_6(3),\ R_8(2).
\]

\section{Formal verification and trust boundary}

The accompanying Lean~4 development is pinned to Lean and mathlib version $4.12.0$.  It
formalizes:
\begin{enumerate}
  \item the exact fixed-width queue semantics and the history-soundness lemma;
  \item the pulse automaton and its exhaustive block-forcing theorem;
  \item both directions of Theorem~\ref{thm:source} for every admissible $\beta,q$;
  \item the fresh-marker construction, fixed-length binary recoding, solvability, and
        primitive terminality of Theorem~\ref{thm:five-pcp};
  \item the four-element source alphabet, nonerasing morphisms, and automatic nonemptiness
        of every GPCP witness, including the explicit $\operatorname{GPCP}^{+}(4)$ form;
  \item the ternary morphism, exact integer matrices, arbitrary-product converse,
        upper-triangularity, nonsingularity, nonzeroness, and rational rank one;
  \item the composed equivalences between tag halting, four-generator GPCP, corrected
        five-pair PCP, and five-matrix mortality.
\end{enumerate}
The concrete matrix label type is $\operatorname{Option}(\texttt{NearyTile})$;
$\texttt{NearyTile}$ has four elements and the option type has five.  Hence the emitted image
contains at most five matrix values, as required by the definition of $\Mort_3(5)$; no
pairwise-distinctness premise is used.  Integer products are cast entrywise to rational
matrices, and the development proves that this cast preserves products and reflects zero.

The source contains no \texttt{sorry}, \texttt{admit}, or custom axiom.  Lean's
\texttt{\#print axioms} command reports only \texttt{propext},
\texttt{Classical.choice}, and \texttt{Quot.sound} for the principal theorems.  Mathlib is
not an additional oracle: its theorem constants carry proof terms rechecked by Lean's kernel.
The trusted base is therefore the Lean kernel and toolchain implementation, together with
Lean's standard classical quotient type theory and the hardware/runtime executing the checker.

The formal artifact does not presently prove Neary's Lemma~9 from a machine model already
known undecidable in mathlib, nor does it formalize the bibliographic antichain reductions.
Thus the novel source and matrix equivalences are machine-checked, while the final
computability-theoretic wrapper imports Neary's peer-reviewed tag-universality theorem.
Formalizing his Table~2 compiler would remove that remaining external mathematical
dependency; adding it as an axiom would merely conceal, rather than reduce, the trust base.

\begin{thebibliography}{10}

\bibitem{BPS21}
P.~C. Bell, I. Potapov, and P. Semukhin.
On the mortality problem: From multiplicative matrix equations to linear recurrence
sequences and beyond.
\emph{Information and Computation}, 281:104736, 2021.
\url{https://doi.org/10.1016/j.ic.2021.104736}.

\bibitem{CHHN14}
J. Cassaigne, V. Halava, T. Harju, and F. Nicolas.
Tighter undecidability bounds for matrix mortality, zero-in-the-corner problems, and more.
arXiv:1404.0644v3, 2014.
\url{https://arxiv.org/abs/1404.0644}.

\bibitem{Dong23}
R. Dong.
Recent advances in algorithmic problems for semigroups.
\emph{ACM SIGLOG News}, 10(4):3--23, 2023.
\url{https://doi.org/10.1145/3636362.3636365}.

\bibitem{HHH07}
V. Halava, T. Harju, and M. Hirvensalo.
Undecidability bounds for integer matrices using Claus instances.
\emph{International Journal of Foundations of Computer Science}, 18(5):931--948, 2007.
\url{https://doi.org/10.1142/S0129054107005066}.

\bibitem{Neary13}
T. Neary.
Undecidability in binary tag systems and the Post correspondence problem for four pairs of
words.
arXiv:1312.6700v1, 2013.  One-version technical report; claim not retained in the refereed
successor.
\url{https://arxiv.org/abs/1312.6700}.

\bibitem{Neary15}
T. Neary.
Undecidability in binary tag systems and the Post correspondence problem for five pairs of
words.
In \emph{STACS 2015}, LIPIcs 30, pages 649--661, 2015.
\url{https://doi.org/10.4230/LIPIcs.STACS.2015.649}.

\bibitem{Nicolas08}
F. Nicolas.
(Generalized) Post correspondence problem and semi-Thue systems.
arXiv:0802.0726v5, 2008.
\url{https://arxiv.org/abs/0802.0726}.

\bibitem{Rote25}
G. Rote.
Probabilistic finite automaton emptiness is undecidable for a fixed automaton.
In \emph{MFCS 2025}, LIPIcs 345, article 86, 2025.
\url{https://doi.org/10.4230/LIPIcs.MFCS.2025.86}.

\end{thebibliography}

\end{document}
